\chapter{Introduction}\label{ch:intro}

\NF{} is a headless geometry kernel based on non-uniform rational B-splines (NURBS). It builds
exact two- and three-dimensional curves, surfaces, trimmed surfaces and closed solids and hands
them to downstream software: optical design tools, finite element (FEM), boundary element
(BEM) and fast multipole (FMM) solvers, mesh generators, slicers for additive manufacturing and
other CAD systems. There is no graphical user interface. Geometry is created from code, in C++
or Python, or from natural language through a local large language model (LLM). It is meshed
with \href{https://gmsh.info}{gmsh} and exchanged through STEP files.

This manual describes \NF{} version \NFVersion. It explains the concepts behind the kernel,
walks through every area of the Python interface with tested examples and closes with a
complete command reference in \cref{app:reference}.

\section{What \NF{} is for}\label{sec:intro-purpose}

Scientific software often needs geometry that is \emph{exact} and \emph{reproducible}: a lens
surface whose sag must be correct to a nanometre, a scatterer whose boundary integral must
converge at the expected rate, or a micro-optical element whose writing paths must follow the
true surface. General-purpose CAD systems are interactive, and their file formats lose
information or need large libraries to read. \NF{} closes this gap with three properties:

\begin{itemize}
  \item \textbf{Exact representations.} Lines, circles, conics, spheres, cylinders, cones and
    tori are represented exactly as rational NURBS. Operations that cannot be exact (fits,
    sweeps, offsets, optical sags, intersection curves) take an explicit tolerance and check
    the result against it.
  \item \textbf{Scriptability.} Every operation is a function call. A model is a short Python
    script that can be versioned, parametrised, optimised and run on a cluster without a
    display.
  \item \textbf{Embeddability.} The C++ kernel has no third-party runtime dependencies. It can be
    linked into a solver as a static library. gmsh and Ollama are used only from the optional
    Python layer.
\end{itemize}

\subsection{Typical users}

\begin{description}[style=nextline,leftmargin=1.5em]
  \item[Optical design] Aspheres, Zernike and XY-polynomial freeform surfaces, exact conics,
    circular apertures as trimmed surfaces, and surface curvature for tolerance analysis
    (\cref{ch:optics}).
  \item[FEM, BEM and FMM solver developers] Closed B-rep solids from primitives and Booleans,
    face labels that become physical groups, exact NURBS transfer to gmsh, local mesh sizes,
    and watertight, outward-oriented surface meshes for boundary element methods
    (\cref{ch:solids,ch:meshing}).
  \item[Two-photon lithography (2PP) and additive manufacturing] Planar sections of solids,
    critical heights at which the layer topology changes, and hatch and contour lines as
    writing paths (\cref{ch:slicing}).
  \item[CAD data exchange] A native, dependency-free STEP (AP214) writer and reader with units,
    solid names and exact analytic faces (\cref{ch:step}).
  \item[Exploratory and automated modelling] An LLM interface that turns requests in natural
    language into calls of whitelisted geometry tools, checks the results against measurements
    and solves dimension constraints (\cref{ch:llm}).
\end{description}

\section{Feature overview}\label{sec:intro-features}

\begin{longtable}{@{}p{0.22\textwidth}p{0.72\textwidth}@{}}
\toprule
\textbf{Area} & \textbf{Capabilities} \\
\midrule
\endhead
NURBS core & Curves and tensor-product surfaces of any degree, rational or not; evaluation and
  derivatives of any order; knot insertion and refinement, knot removal, degree elevation and
  reduction, splitting, joining, Bézier decomposition, clamping, reversal, transposition. \\
Primitives & Points, lines, polylines, exact arcs and circles (also through three points),
  regular polygons, rectangle outlines, helices, conical helices and spirals, bilinear and
  rectangular patches. \\
Construction & Extrusion, revolution, ruled surfaces, loft (skinning), sweeps with one or two
  rails, Coons patches, Gordon surfaces, extrapolation, offsets of curves and surfaces, curve
  and surface fillets, corner rounding. \\
Fitting & Interpolation and least-squares approximation of curves and surfaces, fits to a
  tolerance, end tangents, several parametrisation methods. \\
Optics & Exact conics, even and odd aspheres, Zernike (Fringe and Noll), XY polynomials and
  arbitrary sag functions, fitted to a tolerance. \\
Analysis & Closest points, curvature and torsion, surface curvature, arc length, area,
  curve/curve, curve/surface, curve/plane and surface/surface intersections, planar
  sections, numerical integration. \\
Topology & Trimmed surfaces with exact trim loops, splitting along intersections, B-rep shells
  (sewing, watertightness, orientation), solid primitives, Boolean union, intersection and
  difference, point classification, mass properties, face labels. \\
Slicing & Sections of closed shells into oriented loops, critical heights, extrusion-band
  detection, hatch and contour lines for writing paths. \\
Exchange & Native STEP writer and reader (AP203/AP214/AP242 B-rep), units, solid names. \\
Meshing & Exact transfer to gmsh's OpenCASCADE kernel, physical groups from face labels, local
  mesh sizes, curvature-based sizes, high-order elements, BEM-ready surface meshes. \\
LLM & Ollama tool calling with a whitelist of geometry tools, verification rounds, constraint
  solving, an interactive command-line shell. \\
\bottomrule
\end{longtable}

\section{Design principles}\label{sec:intro-principles}

The following principles shape the whole interface. Knowing them makes the behaviour of
individual functions predictable.

\subsection{Exactness first}
Whatever can be represented exactly is represented exactly. A circle is a rational quadratic
NURBS curve whose points lie on the circle up to floating-point rounding, about $10^{-16}$
relative. Spheres, cylinders and tori are revolved exact profiles. The data passed to gmsh or
written to STEP are the original NURBS (degree, knots, multiplicities, weights), never a
tessellation.

Where exactness is impossible, the operation takes a \code{tolerance} in model units and
\emph{verifies} it. This applies to fitting a sweep, offsetting a curved surface, tracing an
intersection curve, or fitting an aspheric sag. The result is checked at dense check points
between the data that defined it, and the construction is refined until the check passes. The
tolerance model is described in \cref{sec:concepts-tolerances}.

\subsection{Value semantics and immutability}
Geometry objects are immutable values. An operation never modifies its input. It returns a new
object, for example \code{curve.insert_knot(u)}, \code{surface.transformed(T)} or
\code{nf.boolean(a, b, op)}. A curve stored in a list cannot change behind your back, and
objects can be shared freely between threads.

\subsection{Validated invariants}
Constructors check their input: the knot vector length, monotonicity and multiplicities,
positive weights, net sizes, non-degenerate axes. Invalid input raises
\code{nf.GeometryError}, a subclass of Python's \code{ValueError}, with a message that names
the violated condition. An object that exists is therefore always valid.

\subsection{Headless and scriptable}
Every operation is reachable from Python. The LLM layer only composes existing API calls, and
the language model never executes code (\cref{ch:llm}).

\subsection{Portable, dependency-free kernel}
The kernel is written in standard C++20 and needs only the standard library and the system
thread library. It is tested on every change with GCC (AlmaLinux 10), MSVC (Windows 11) and
Apple Clang (macOS, Apple Silicon), and additionally under ThreadSanitizer.

\section{Architecture}\label{sec:intro-architecture}

\NF{} consists of three layers. Dependencies point strictly downwards:

\begin{lstlisting}[style=text]
 +--------------------------------------------------------------+
 |  Python package  nurbforge                                   |
 |   +- llm/   Ollama client, whitelisted tools, agent, CLI     |
 |   +- mesh/  exact NURBS transfer to gmsh (OCC), meshing      |
 |   +- patterns.py, __init__ (re-exports of _core)             |
 +--------------------------------------------------------------+
 |  nurbforge._core  (pybind11 bindings)                        |
 +--------------------------------------------------------------+
 |  C++ kernel  libnurbforge_core  (namespace nf)               |
 |   io/        native STEP writer and reader                   |
 |   topology/  trimmed surfaces, shells, Booleans, solids      |
 |   optics/    conics, aspheres, Zernike, XY polynomials       |
 |   modeling/  construct, loft, sweep, Coons/Gordon, extend,   |
 |              offset, fillet, patterns                        |
 |   fitting/   interpolation, least-squares approximation      |
 |   analysis/  closest point, properties, intersections,       |
 |              sections, hatching                              |
 |   geometry/  transforms, primitives                          |
 |   nurbs/     knot vectors, basis functions, curves, surfaces |
 |   core/      Vec3/Vec4, tolerances, errors                   |
 +--------------------------------------------------------------+
\end{lstlisting}

\begin{description}[style=nextline,leftmargin=1.5em]
  \item[C++ kernel] All geometric algorithms. It is a static library
    (\code{Nurbforge::core}) that can be used without Python (\cref{ch:cpp}).
  \item[\code{nurbforge._core}] A thin pybind11 module that exposes the kernel types and
    functions one to one. All of its names are re-exported by the top-level package, so
    you write \code{nf.make_box(...)}, never \code{nf._core.make_box(...)}.
  \item[\code{nurbforge.mesh}] Pure Python on top of the gmsh API (optional dependency).
  \item[\code{nurbforge.llm}] Pure Python on top of the Ollama HTTP API (standard library
    only).
\end{description}

\section{Platforms and versions}\label{sec:intro-platforms}

\begin{tabularx}{\textwidth}{@{}lX@{}}
\toprule
\textbf{Platform} & \textbf{Support} \\
\midrule
Linux x86\_64 & Binary wheels \code{manylinux_2_28} (AlmaLinux/RHEL 8, 9, 10 and compatible);
  source builds with GCC 14 on AlmaLinux 10 are tested in CI. \\
Windows 11 x64 & Binary wheels \code{win_amd64}; source builds with Visual Studio 2022/2026. \\
macOS 14 or newer & Binary wheels for Apple Silicon (\code{arm64}) and Intel (\code{x86_64});
  source builds with Apple Clang. \\
Python & CPython 3.10, 3.11, 3.12, 3.13 and 3.14. \\
\bottomrule
\end{tabularx}

\paragraph{Versioning and API stability.}
\NF{} follows semantic versioning. As long as the major version is 0, the API is \emph{not yet
stable}: minor releases (0.8 $\to$ 0.9) may change names or signatures. Every change is listed
in \code{CHANGELOG.md}. The version of an installed package is available as
\code{nf.__version__}, and in C++ as the macro \code{NURBFORGE_VERSION}. Planned work is listed in
the roadmap (\code{docs/roadmap.md}).

\section{How to read this manual}\label{sec:intro-reading}

\begin{itemize}
  \item \Cref{ch:install,ch:quickstart} get you from installation to a first meshed solid in a
    few minutes.
  \item \Cref{ch:concepts} explains NURBS, parameter domains, tolerances and error handling.
    Read it before relying on numerical results.
  \item \Cref{ch:curves,ch:surfaces,ch:fitting,ch:optics,ch:intersections,ch:trimming,%
ch:solids,ch:slicing,ch:transforms} cover the modelling functions area by area.
  \item \Cref{ch:step,ch:meshing,ch:llm} describe data exchange, meshing and the LLM interface,
    and \cref{ch:cpp} the C++ interface.
  \item \Cref{ch:troubleshooting} collects common problems and their solutions.
  \item The appendices are references. \Cref{app:reference} lists every Python class,
    method and function with its signature and description, generated from the installed
    package. \Cref{app:llmtools} lists the LLM tools, \cref{app:cli} the command-line
    interface, and \cref{app:tolerances} every tolerance and default.
\end{itemize}

\subsection{Typographic conventions}

\begin{tabularx}{\textwidth}{@{}lX@{}}
\toprule
\textbf{Notation} & \textbf{Meaning} \\
\midrule
\code{make_box} & Names of functions, classes, parameters and literal code. \\
\apiref{make_box} & A name linked to its entry in the command reference (in the PDF, click it). \\
\code{nf.} & The package imported as \code{import nurbforge as nf}, used in all examples. \\
$C(u)$, $S(u,v)$ & Mathematical curves and surfaces with their parameters. \\
\code{examples/03_quickstart.py} & A complete script shipped in \code{docs/manual/examples/}. \\
\bottomrule
\end{tabularx}

Code listings with a file name in the title are complete scripts. The test suite runs every one
of them, so they work with the version of \NF{} this manual describes. Numbers quoted in the
text come from running these scripts. Short listings without a title are fragments that assume
the objects of the surrounding text.

\begin{note}[Boxes]
Green boxes contain background notes, blue boxes practical tips, and orange boxes warnings
about pitfalls that lead to wrong results or errors.
\end{note}

Unless stated otherwise:
\begin{itemize}
  \item angles are in \textbf{radians} (the LLM tools use degrees; \cref{ch:llm});
  \item lengths are in arbitrary but consistent \textbf{model units}; \NF{} attaches a unit only
    when writing or reading STEP files (\cref{sec:concepts-units});
  \item curve and surface parameters run over $[0,1]$ for everything the constructors create.
\end{itemize}

\section{Getting help}\label{sec:intro-help}

\begin{itemize}
  \item Source code, issue tracker and releases: \url{https://github.com/PhotonPower/Nurbforge}.
  \item Online documentation (tutorial, API reference, C++ reference generated by Doxygen):
    \url{https://photonpower.github.io/Nurbforge/}.
  \item In Python, \code{help(nf.make_box)} shows the signature and docstring of any function.
  \item Bug reports should contain the \NF{} version (\code{nf.__version__}), the platform and a
    minimal script that reproduces the problem. For geometry problems, a STEP file written with
    \code{nf.write_step} is often the quickest way to share the input.
\end{itemize}

\NF{} is released under the MIT licence.
